\begin{thebibliography}{99}
\bibitem[Bekolay et~al.(2014)]{bekolay2014}
Bekolay, T., Bergstra, J., Hunsberger, E., DeWolf, T., Stewart, T.~C., Rasmussen, D., Choo, X., Voelker, A.~R., and Eliasmith, C. (2014).
Nengo: A Python tool for building large-scale functional brain models.
\textit{Frontiers in Neuroinformatics}, 7, 48.
\url{https://doi.org/10.3389/fninf.2013.00048}.
\bibitem[Crawford et~al.(2016)]{crawford2016}
Crawford, E., Gingerich, M., and Eliasmith, C. (2016).
Biologically plausible, human-scale knowledge representation.
\textit{Cognitive Science}, 40, 782--821.
\url{https://doi.org/10.1111/cogs.12261}.
\bibitem[Dell et~al.(1997)]{dell1997}
Dell, G.~S., Schwartz, M.~F., Martin, N., Saffran, E.~M., and Gagnon, D.~A. (1997).
Lexical access in aphasic and nonaphasic speakers.
\textit{Psychological Review}, 104(4), 801--838.
\url{https://doi.org/10.1037/0033-295X.104.4.801}.
\bibitem[Eliasmith(2013)]{eliasmith2013}
Eliasmith, C. (2013).
\textit{How to Build a Brain: A Neural Architecture for Biological Cognition}.
Oxford University Press.
\bibitem[Eliasmith and Anderson(2003)]{eliasmith2003}
Eliasmith, C., and Anderson, C.~H. (2003).
\textit{Neural Engineering: Computation, Representation, and Dynamics in Neurobiological Systems}.
MIT Press.
\bibitem[Eliasmith and Martens(2011)]{eliasmithmartens2011}
Eliasmith, C., and Martens, J. (2011).
Normalization for probabilistic inference with neurons.
\textit{Biological Cybernetics}, 104(4--5), 251--262.
\url{https://doi.org/10.1007/s00422-011-0433-y}.
\bibitem[Eliasmith et~al.(2012)]{eliasmith2012}
Eliasmith, C., Stewart, T.~C., Choo, X., Bekolay, T., DeWolf, T., Tang, Y., and Rasmussen, D. (2012).
A large-scale model of the functioning brain.
\textit{Science}, 338(6111), 1202--1205.
\url{https://doi.org/10.1126/science.1225266}.
\bibitem[Kröger(2023)]{kroeger2023}
Kröger, B.~J. (2023).
The NEF-SPA approach as a framework for developing a neurobiologically inspired spiking neural network model for speech production.
\textit{Journal of Integrative Neuroscience}, 22(5), 124.
\url{https://doi.org/10.31083/j.jin2205124}.
\bibitem[Kröger(2026a)]{kroegerArticulation2026}
Kröger, B.~J. (2026a).
From syllable representation to articulation using a spiking neural model approach.
Companion working manuscript, 5 November 2026; not yet submitted.
\bibitem[Kröger(2026b)]{kroegerAcoustics2026}
Kröger, B.~J. (2026b).
Vocal-tract acoustics, aerodynamics, and overlap-and-add synthesis in DYNARTmo.
\textit{arXiv}, 2610.24490.
\url{https://arxiv.org/abs/2610.24490}.
\bibitem[Kröger et~al.(2016)]{kroeger2016}
Kröger, B.~J., Crawford, E., Bekolay, T., and Eliasmith, C. (2016).
Modeling interactions between speech production and perception: Speech error detection at semantic and phonological levels and the inner speech loop.
\textit{Frontiers in Computational Neuroscience}, 10, 51.
\url{https://doi.org/10.3389/fncom.2016.00051}.
\bibitem[Levelt et~al.(1999)]{levelt1999}
Levelt, W.~J.~M., Roelofs, A., and Meyer, A.~S. (1999).
A theory of lexical access in speech production.
\textit{Behavioral and Brain Sciences}, 22(1), 1--38.
\url{https://doi.org/10.1017/S0140525X99001776}.
\bibitem[Ma et~al.(2006)]{ma2006}
Ma, W.~J., Beck, J.~M., Latham, P.~E., and Pouget, A. (2006).
Bayesian inference with probabilistic population codes.
\textit{Nature Neuroscience}, 9(11), 1432--1438.
\url{https://doi.org/10.1038/nn1790}.
\bibitem[NengoSPA developers(n.d.)]{nengospaAM}
NengoSPA developers (n.d.).
Associative memory.
\textit{NengoSPA documentation}.
\url{https://www.nengo.ai/nengo-spa/examples/associative-memory.html} (accessed 6 November 2026).
\bibitem[Stewart et~al.(2011)]{stewart2011}
Stewart, T.~C., Tang, Y., and Eliasmith, C. (2011).
A biologically realistic cleanup memory: Autoassociation in spiking neurons.
\textit{Cognitive Systems Research}, 12(2), 84--92.
\url{https://doi.org/10.1016/j.cogsys.2010.06.006}.
\bibitem[Stille et~al.(2020)]{stille2020}
Stille, C.~M., Bekolay, T., Blouw, P., and Kröger, B.~J. (2020).
Modeling the mental lexicon as part of long-term and working memory and simulating lexical access in a naming task including semantic and phonological cues.
\textit{Frontiers in Psychology}, 11, 1594.
\url{https://doi.org/10.3389/fpsyg.2020.01594}.
\end{thebibliography}
